The work below is an editorial verification of three precise statements in the
notebook.  It makes no claim of a new general theorem or of a solution to an
unresolved problem in the problem index.  Results proved here are explicitly
marked; historical identities retain their attribution.  The decision about
research development of the class C and D items is recorded in the overview
and report.

\subsection{The grading check for Atiyah's de Rham example}
\label{verification:grading}

On \GSpage{P000373}, the notes use \(d+\delta\) on differential forms as an
example in the discussion of spectral flow and an obstruction to an invariant
metric.  The full de Rham operator needs to be distinguished from the odd
signature operator.  The following elementary check concerns the full operator
and does not change the general spectral-flow/index theorem recorded in the
preceding pages.
The general index identification used below is
\citet[\S7, Theorem 7.4 and pp.94--95]{atiyahpatodisinger1976spectral3}.

\begin{proposition}[Vanishing for a graded loop]\ours{}
Let \(D_t\), \(0\leq t\leq1\), be a smooth family of self-adjoint first-order
elliptic operators on a fixed Hermitian vector bundle over a closed smooth
manifold.  Regard them as operators on a fixed Hilbert space with the same
Sobolev domain.  Suppose \(D_1=UD_0U^{-1}\), where \(U\) is the unitary gluing
operator, and suppose a fixed unitary involution \(\Gamma\) preserves the
domains, commutes with \(U\), and satisfies
\[
\Gamma D_t\Gamma=-D_t\qquad(0\leq t\leq1).
\]
Then the spectral flow of the resulting closed family is zero.  In particular,
the full de Rham family \(d+\delta_{g_t}\), glued by the action of a
diffeomorphism and identified with a fixed Hilbert space in a
degree-preserving way, has zero spectral flow.
\end{proposition}

\begin{proof}
Spectral flow is unchanged by a unitary conjugacy of the closed family, and
replacing every operator by its negative reverses the oriented crossing
count.  These properties follow first for a family with isolated regular
crossings, and then for arbitrary families by perturbation and homotopy
invariance.  For a closed family the dimensions of the endpoint kernels
agree, so no endpoint correction survives.  Consequently,
\[
\operatorname{sf}(D_t)=\operatorname{sf}(\Gamma D_t\Gamma)
=\operatorname{sf}(-D_t)=-\operatorname{sf}(D_t).
\]
The spectral flow is an integer, so it is zero.

For the last assertion, let \(\Gamma\) act by \((-1)^k\) on forms of degree
\(k\).  Both \(d\) and \(\delta_{g_t}\) reverse this parity, and a
diffeomorphism preserves degree.  The positive changes of density and of the
fibre metric used to identify the varying \(L^2\) spaces are block diagonal
by degree and hence commute with \(\Gamma\).  All the required identities
therefore hold after those identifications.
\end{proof}

\paragraph{Independent check and scope.}
For geometric bundle gluing, the same conclusion can be checked by the
suspended operator.  On the mapping torus use a unitary connection along the circle which preserves
\(\Gamma\), and put \(P=\nabla_t+D_t\).  Integration by parts gives
\(P^*=-\nabla_t+D_t\), while
\(\Gamma P\Gamma=\nabla_t-D_t=-P^*\).  Unitary conjugation preserves index
and the adjoint changes its sign, whence \(\operatorname{ind}P=0\).  This is
an independent algebraic check of the vanishing once the historical
spectral-flow/index identification is used.  It explains why an unspecified
reference to \(d+\delta\) cannot, by itself, supply a nonzero obstruction.
It does not assert that every natural self-adjoint elliptic family has zero
spectral flow, nor does it identify a signature grading that the notes do not
specify.

\subsection{An exact finite audit of Dyson's coefficient formula}
\label{verification:dyson}

The product printed on \GSpage{P000445} needs the normalization
\[
\Delta(q)=q\prod_{m=1}^{\infty}(1-q^m)^{24}
=\sum_{n=1}^{\infty}\tau(n)q^n.
\]
The leading \(q\) is already necessary on comparison of constant terms:
the unshifted product has constant coefficient one and the displayed series
has constant coefficient zero.  This is a mathematical emendation of the
notebook formula, not an inference about an invisible pen stroke.
The historical all-coefficient formula is
\citet[pp.636--637, equations (1)--(3)]{dyson1972missed}.

\begin{proposition}[Finite coefficient agreement]\oursComputed{}
For every integer \(1\leq n\leq100\), the coefficient of \(q^n\) in the
corrected product agrees exactly with
\[
\sum_{\substack{
 (a,b,c,d,e)\equiv(1,2,3,4,0)\pmod5\\
 a+b+c+d+e=0\\a^2+b^2+c^2+d^2+e^2=10n}}
\frac{\displaystyle\prod_{1\leq i<j\leq5}(v_i-v_j)}{1!2!3!4!},
\qquad v=(a,b,c,d,e).
\]
This assertion reports a finite exact computation and is not a proof of the
all-\(n\) identity.
\end{proposition}

\paragraph{Algorithms and exhaustiveness.}
The file \texttt{code/check\_dyson.py} first multiplies the finite binomial
expansions of \((1-q^m)^{24}\), retaining terms only through the requested
degree.  Factors with larger \(m\) cannot change those coefficients.  A second
algorithm computes the coefficients \(a_n\) of the unshifted product from
its formal logarithmic derivative:
\[
a_0=1,\qquad
n a_n=-24\sum_{j=1}^{n}\sigma_1(j)a_{n-j},\qquad
\tau(n)=a_{n-1}.
\]
The third algorithm enumerates the first four lattice coordinates with the
specified residues and sets the fifth equal to the negative of their sum.
For a run through \(N\), every coordinate is bounded in absolute value by
\(\lfloor\sqrt{10N}\rfloor\); the squared-length condition then selects
the coefficient.  Thus the enumeration omits no qualifying tuple.  All
arithmetic is integral, including explicit divisibility checks for the
denominator \(1!2!3!4!=288\).

\begin{center}
\begin{tabular}{r r r r}
\toprule
\(n\) & \(\tau(n)\) & \(n\) & \(\tau(n)\)\\
\midrule
1&1&7&-16744\\
2&-24&8&84480\\
3&252&9&-113643\\
4&-1472&10&-115920\\
5&4830&11&534612\\
6&-6048&12&-370944\\
\bottomrule
\end{tabular}
\end{center}

\paragraph{Second check and limitation.}
The complete comparison was run separately with \(N=32\) and \(N=100\);
all three arrays agreed in both runs.  The two output files in
\texttt{audit/} record every checked coefficient and the number of contributing
lattice tuples.  Agreement also checks the residue ordering and the sign of
the ten-factor Vandermonde as copied.  The verification supports the
normalization correction; it does not convert finite evidence into an
identity for every integer, and it has no bearing on the separate question
about the exponent \(26\).
An independent review also checked the labelled residue classes, the absence
of a factorial permutation divisor, the enumeration bound, and the recurrence;
it found no error.

\subsection{Recovering the eighty-million lattice bound}
\label{verification:mass}

The numerical remark on \GSpage{P000398} becomes identifiable after restoring
positive definiteness, integrality and evenness.  For positive-definite even
unimodular lattices of rank \(n=8k>0\), the classical mass formula is
\[
m_n=\sum_{[L]}\frac{1}{|\operatorname{Aut}(L)|}
=\frac{|B_{n/2}|}{n}\prod_{j=1}^{n/2-1}\frac{|B_{2j}|}{4j},
\]
where \(t/(e^t-1)=\sum_{j\geq0}B_jt^j/j!\).
The formula and comparison bound are recorded by
\citet[\S\S2, 6]{king2003mass}.

\begin{proposition}[Exact mass evaluation]\oursComputed{}
Evaluation in rational arithmetic gives
\[
m_{32}=40309232.93560208\ldots,\qquad
\lceil 2m_{32}\rceil=80618466.
\]
Thus the mass formula alone proves that there are at least \(80618466\)
isometry classes in this genus.  It does not give their exact number.
\end{proposition}

\begin{proof}
The program \texttt{code/check\_mass.py} evaluates the displayed historical
formula using exact fractions; its full numerator and denominator are retained
in \texttt{audit/mass\_check.json}.  Every positive-rank lattice admits the two
distinct isometries \(I\) and \(-I\), so each summand of the mass is at most
\(1/2\).  The number of summands is therefore at least \(2m_{32}\), and
integrality gives the ceiling.
\end{proof}

\paragraph{Independent arithmetic check.}
The Bernoulli numbers were computed both from
\(\sum_{j=0}^{r}\binom{r+1}{j}B_j=0\) for \(r\geq1\), with \(B_0=1\),
and from
\(B_r=\sum_{j=0}^{r}(-1)^j j!S(r,j)/(j+1)\), where \(S(r,j)\) are Stirling
numbers of the second kind.  All values through \(B_{30}\) agreed.
As normalization checks, the same computation gives
\(m_8=1/696729600\) and
\[
m_{16}=\frac{1}{2(696729600)^2}+\frac{1}{2^{15}16!},
\]
the reciprocal-automorphism sum for \(E_8\oplus E_8\) and \(D_{16}^{+}\).
All checks passed.  This recovers the notebook's rough magnitude; the later
stronger lower bound belongs to the cited literature and is not recomputed
here.
